\section{Study Area}
\label{sec:study_area}

The \emph{Stagnone di Marsala} is located on the north-western coast of
Sicily, sheltered from the western Sicilian channel by Isola Grande, an
elongated barrier island running parallel to the mainland coast
(Figure~\ref{fig:study_area}). The enclosed basin extends approximately 11~km
in the north--south direction and 2.5~km east--west, enclosing a total surface
area of roughly 2200~ha \citep{DeMarchis2012}. Two geomorphologically distinct
sub-basins are separated by a \emph{Posidonia oceanica} barrier reef
\citep{LaLoggia2004}: the northern sub-basin, with a mean depth of
approximately 1~m and extensive intertidal flats, and the deeper southern
sub-basin, with a mean depth of approximately 2~m and a maximum of about 3~m
near the main inlet. Three islands punctuate the interior, the Phoenician
archaeological site of Mothia in the south-central basin and the smaller
islands of Santa Maria and Scola in the central reaches
\citep{CiraoloDeMarchis2009}. Salt pans (\emph{saline}) occupy both margins of
the lagoon, with extensive historic saltworks on the western barrier island
(Isola Grande) and additional pans along the eastern mainland shore in the
area facing Mothia \citep{Basso2008}. They constitute a distinctive shallow
intertidal habitat within the lagoon system.

Hydrodynamic exchange with the open sea occurs through two morphologically
distinct inlets. The northern inlet (approximately 400~m wide, 0.30--0.40~m
deep) is a highly restricted passage whose tidal prism is severely limited by
its shallow sill. A historically dredged channel (20~m wide, 1~m deep)
partially compensates this restriction \citep{CiraoloDeMarchis2009}. The wide
southern inlet (approximately 2900~m wide, 1.0--1.5~m deep) provides the
primary pathway for mass transport and water mixing between the lagoon and the
open sea \citep{CiraoloDeMarchis2009}. The tidal regime is microtidal, with
free-surface amplitudes of approximately 0.30~m dominated by semi-diurnal
harmonics \citep{LaLoggia2004}. In the absence of riverine input, weak tidal
exchange combined with high summer evaporation drives a pronounced salinity
gradient. Annual lagoon-wide salinity spans 33--46~psu \citep{Basso2008}, with
peak summer values approaching 48~psu in the more confined northern sub-basin,
contrasting with near-marine conditions near the southern inlet
\citep{Mancuso2023}. Wind forcing dominates east--west circulation within the
basin, while tidal exchange controls the north--south component
\citep{CiraoloDeMarchis2009, DeMarchis2012}.

The seagrass mosaic covers the majority of the lagoon floor.
\emph{Posidonia oceanica} occupies the central and southern portions of the
basin, forming both continuous meadows and characteristic atoll patterns
(10--20~m in diameter) that grade into reef formations oriented along the
north--south axis. \emph{Cymodocea nodosa} prevails in the shallower northern
reaches, with \emph{Caulerpa prolifera} replacing \emph{P.\ oceanica} locally
where hydrodynamic and sediment conditions deteriorate
\citep{LaLoggia2004, CiraoloDeMarchis2009}. In Mediterranean meadows of this
type, \emph{P.\ oceanica} areal densities typically range from 500 to
1200 plants~\si{\per\metre\squared} \citep{Ciraolo2009a}. The spatial
distribution of submerged vegetation was mapped from satellite imagery by
\citet{CiraoloMaltese2006}. The roughness parametrisation in this study uses
the updated classification by \citet{Maltese2025} as the base map, applied to
the 2023 epoch (Section~\ref{sec:methods_roughness}). The lagoon is designated
both as a \emph{Riserva Naturale Orientata} and as a Natura 2000 site,
reflecting the conservation importance of these meadows, which have
nevertheless declined by approximately 75\% in extent over the past two
decades \citep{Maltese2025}.

The observational dataset used in this study was collected during July 2025
(Figure~\ref{fig:study_area}). Water level was recorded at three stations,
BocaNord (BN) at the northern inlet, BocaSud (BS) near the southern inlet, and
AltaVilaEst (AE) on the eastern flank of the central basin. Wind speed and
direction and atmospheric pressure were measured at two stations, AE (which
also carries a water-level gauge and a water temperature sensor) and Mulino, a
dedicated meteorological station located near the salt pans in the
central-northern area. Water levels at BN, BS, and AE serve as the primary
validation targets for the hydrodynamic model
(Section~\ref{sec:methods_validation}). Open-sea boundary conditions are
prescribed from CMEMS Mediterranean physical reanalysis products
(Section~\ref{sec:methods_forcing}). Offshore model performance is assessed
against the tidal gauge on Marettimo island (JRC TAD 658), approximately 15~km
to the west, which provides a free-surface reference uninfluenced by lagoonal
dynamics. GPS-tracked surface drifters were released in the lagoon interior
over a sub-period coinciding with the simulation window, and their
trajectories form the basis of the Lagrangian validation
(Section~\ref{sec:methods_validation}).

%% TODO Figure 1 (study area map) not yet produced. Needs: lagoon bathymetry,
%% the two inlets, the three islands, the four in-situ stations (BN, BS, AE,
%% Mulino), Marettimo inset, and the drifter release points. Replace the
%% placeholder box below with \includegraphics once the figure exists.
\begin{figure}[H]
  \centering
  \fbox{\parbox[c][6cm][c]{0.9\textwidth}{\centering\textsf{[Figure~1
  placeholder] Study area map: lagoon bathymetry, northern and southern
  inlets, Mothia/Santa Maria/Scola, in-situ stations BN, BS, AE and Mulino,
  drifter release points, Marettimo inset.}}}
  \caption{The Stagnone di Marsala lagoon on the north-western coast of
  Sicily, showing bathymetry, the two inlets, the interior islands, and the
  July 2025 in-situ station network.}
  \label{fig:study_area}
\end{figure}
